\documentclass[10pt,a4paper]{amsart}
\usepackage[margin=19mm]{geometry}
\usepackage{amsmath,amssymb,mathtools}
\usepackage[scr=rsfs,cal=euler]{mathalfa}
\usepackage{xfrac}
\usepackage[unicode,hidelinks,bookmarks=false]{hyperref}
\newcommand{\ie}{i.e.\ }
\newcommand{\cf}{cf.\ }
\newcommand{\resp}{respectively\ }
\newcommand{\Subsection}{\subsection}
\providecommand{\nobreakdash}{\nobreak\hbox{-}}
\setlength{\emergencystretch}{2em}
\multlinegap0pt
\usepackage{bm}
\usepackage{bbm}
\usepackage{dsfont}
\usepackage{xspace}
\usepackage{cancel}
\usepackage{mathtools}
\usepackage{amsthm}
\makeatletter
\@ifundefined{lemma}{\newtheorem{lemma}{Lemma}}{}
\makeatother
\title[Travel Bans vs. Other Disease Mitigation Measures: A Mathema]{Travel Bans vs. Other Disease Mitigation Measures: A Mathematical Analysis}
\author{Christian Borgs, Karissa Huang,, Geng Zhao}
\date{Source: arXiv:2510.08895v2 (CC BY 4.0)}
\begin{document}
\maketitle

\noindent\emph{Source and licence.} Travel Bans vs. Other Disease Mitigation Measures: A Mathematical Analysis. Version arXiv:2510.08895v2 (CC BY 4.0), distributed under the Creative Commons Attribution 4.0 International licence, as recorded in the arXiv licence metadata for this exact version and shown on the abstract page https://arxiv.org/abs/2510.08895v2. Full rights notice: licenses/arxiv-2510.08895-CC-BY-4.0.txt.\par\smallskip

\noindent\textit{Editorial scope.} Counted unit: the Lemma whose source label is \texttt{lem:matrix\_eigenvalues} in the article (source0.tex lines 1477--1495, source label \texttt{lem:matrix\_eigenvalues}), delivered together with the article's own complete original proof of it (source0.tex lines 1497--1576, one \texttt{proof} environment of 80 lines). No mathematics is written, repaired or re-derived: statement and proof are the article's own text line for line. Required context retained uncounted: M0-def (lines 1465--1470). Every definition, display and label of those lines is retained unchanged. Omitted from this package and not redistributed: 1--1464, 1471--1476, 1496--1496, 1577--2045. Nothing inside a retained excerpt is omitted or altered: every retained body is delivered exactly as the article's lines are written, so no omission receipt of any kind accompanies this package. Citations: the retained excerpts carry no citation command at all, so no bibliography entry of the article is transcribed and no reference is redistributed. Reference adaptation: none was needed and none was made; every reference inside the delivered unit resolves inside it, so no unresolved reference or citation is produced. Printed slips kept literal: no printed word, symbol, hypothesis or number of the counted statement or of its proof was altered. Layout and class: the article's own class is \texttt{article} with packages including fullpage, setspace, graphicx, subcaption, amssymb, amsfonts, amsmath, amsthm, latexsym, epsfig, bm, bbm, mathtools, xspace; the wrapper is the lane's pinned \texttt{amsart} editorial wrapper at 10pt on A4, which restates the theorem-like environments the retained text needs and adds the article's own macro definitions that the retained text needs (none), with extra wrapper packages (bm, bbm, dsfont, xspace, cancel, mathtools). The delivered numbering is the unit's own. Assets: no retained span contains a figure environment, a graphics-inclusion command, a PDF-page inclusion, a table or a TikZ picture; no image file is copied into this package, the package declares no asset dependency and no image file is redistributed with it. Licence: version arXiv:2510.08895v2 of this article is distributed under the Creative Commons Attribution 4.0 International licence, as recorded in the arXiv licence metadata for this exact version and shown on the abstract page https://arxiv.org/abs/2510.08895v2. A full rights notice is delivered at licenses/arxiv-2510.08895-CC-BY-4.0.txt. Characters: every retained excerpt is pure ASCII, so the delivered engine, XeLaTeX with its default Latin Modern text font, reports no missing character.
\begin{equation}\label{M0-def}
(M_0 \mathbf v)_{ij}^{AB}=
 -\left(\gamma+\rho_H\delta_{AT}+\rho_H\delta_{BT}+\beta \delta_{ij}^{AB}\right)v_{ij}^{AB}
+\rho_H\left(\delta_{AH}v_{ij}^{T B}+\delta_{BH}v_{ij}^{AT}\right)
+\beta c \,\delta_{BH}\sum_{k,C}\delta_{ik}^{AC}v_{ki}^{CA}
\end{equation}

\begin{lemma}
\label{lem:matrix_eigenvalues}
The matrix $M_0$ has 16 real-valued eigenvalues, with   $\lambda=(c-1)\beta-\gamma$ being an eigenvalue of algebraic and geometric multiplicity $2$, and all other eigenvalues being $-\gamma$ or smaller.  The projection onto the eigenspace of $\lambda$ is the two-dimensional projection
  $$
  (P_\lambda\mathbf v)=
  \left(\mathbf e_{11}^{HH}+\frac c{c-1}\mathbf e_{12}^{HH}\right)v_{11}^{HH}+
 \left(\mathbf e_{22}^{HH}+\frac c{c-1}\mathbf e_{21}^{HH}\right)v_{22}^{HH},
  $$
  where  we use $\mathbf e_{ij}^{AB}$ to denote the coordinate vector for the coordinate 
\(
\renewcommand{\arraystretch}{0.7}
\left[
\begin{array}{@{}c@{}c@{}c@{}c@{}}
A & B\\
i & j
\end{array}
\right]
\).
\end{lemma}

\begin{proof}
As we will see, after a suitable ordering of the indices, the matrix $M_0$ is lower triangular, allowing us to read off the eigenvalues and their multiplicity from the diagonal.  Assuming this for the moment, we start by computing the diagonal of the matrix $M_0$,
noting that the last term in \eqref{M0-def} gives a contribution to the diagonal if and only if 
$B=H$,
$C=A=B$, $i=j=k$ and
 $\delta_{ik}^{AC}=1$, showing that the only positive contribution to the diagonal are the two entries
$$
\delta_{ij}\delta^{AH}\delta^{BH} (c\beta-\gamma-\beta)=\lambda \delta_{ij}\delta^{AH}\delta^{BH} 
$$
with $i=j=1$ or $i=j=2$, while all others are smaller or equal than $-\gamma$.  This gives the statement about the eigenvalues of $M_0$.

Next we note that 
the restriction of  the matrix $M_0$ 
to the $4$-dimensional subspace of $\mathbb R^{16}$ spanned by 
  $\mathbf e_{11}^{HH}$,   $\mathbf e_{12}^{HH}$,   $\mathbf e_{22}^{HH}$ and  $\mathbf e_{21}^{HH}$ is equal to
  $$
\begin{bmatrix}
    \lambda&0&0&0\\
    c\beta&-\gamma&0&0\\
    0&0&\lambda&0\\
    0&0&c\beta&-\gamma
\end{bmatrix}.
  $$
Diagonalizing this matrix one immediately gets the statement about the projection to the eigenspace of $\lambda$.

We are left with is a proof of the lower triangular structure. 
To this end, we write down the edges $\alpha\to\alpha'$ for which $M_{\alpha',\alpha}\neq 0$ and $\alpha\neq\alpha'$, giving
\begin{align}
\renewcommand{\arraystretch}{0.7}
\left[
\begin{array}{@{}c@{}c@{}c@{}c@{}}   T&T \\ i&i \end{array}\right]
  &\to
\left[
\begin{array}{@{}c@{}c@{}c@{}c@{}}  T&H \\ i&i \end{array}\right],\quad
\left[
\begin{array}{@{}c@{}c@{}c@{}c@{}} T&H \\ i&\bar i \end{array}\right]
     \\
\left[
\begin{array}{@{}c@{}c@{}c@{}c@{}}  H&T \\ \bar i& i \end{array}\right]
  &\to
\left[
\begin{array}{@{}c@{}c@{}c@{}c@{}} T&H \\ i&i \end{array}\right],\quad
\left[
\begin{array}{@{}c@{}c@{}c@{}c@{}}  T&H \\ i&\bar i \end{array}\right]  
  \\
\left[
\begin{array}{@{}c@{}c@{}c@{}c@{}} T&H \\ \bar i& i \end{array}\right]
  &\to
\left[
\begin{array}{@{}c@{}c@{}c@{}c@{}} H&H \\ i&i \end{array}\right],\quad
\left[
\begin{array}{@{}c@{}c@{}c@{}c@{}} H&H \\ i&\bar i \end{array}\right]  
  \\
\left[
\begin{array}{@{}c@{}c@{}c@{}c@{}}  H&H \\ i& i \end{array}\right]
  &\to
\left[
\begin{array}{@{}c@{}c@{}c@{}c@{}}  H&H \\ i&\bar i \end{array}\right]  
\end{align}
where the last line only has one edge since the other is a contribution to the diagonal. If we introduce a partial order by 
\[
\renewcommand{\arraystretch}{0.7}
(TT) > (HT) > (TH) > (HH)
\quad\text{and}\quad
\left[
\begin{array}{@{}c@{}c@{}c@{}c@{}}
H & H\\
i & i
\end{array}\right]
>
\left[
\begin{array}{@{}c@{}c@{}c@{}c@{}}
H & H\\
i & \bar i
\end{array}
\right]
\]
we see that 
all transition corresponding to the last term in \eqref{M0-def} correspond to matrix elements $M_{\alpha',\alpha}$ where $\alpha'<\alpha$.  This clearly also holds for the second term, showing that all off-diagonal terms of $M_0$ respect this order.  Extending the above order to a total order (e.g., by defining a DAG and using topological sorting) this proves that $M_0$ is lower triangular.
\end{proof}

\end{document}
